\documentclass[11pt]{article}
\usepackage[a4paper,margin=25mm]{geometry}
\usepackage{amsmath,amssymb,amsthm,mathtools}
\usepackage[colorlinks=true,urlcolor=blue,citecolor=blue]{hyperref}
\newtheorem{theorem}{Theorem}
\newtheorem{lemma}[theorem]{Lemma}
\DeclareMathOperator{\per}{per}
\DeclareMathOperator{\sgn}{sgn}
\newcommand{\Fp}{\mathbb F_p}
\newcommand{\Zp}{\mathbb Z_{(p)}}
\title{Restricted reciprocal-permutation sums\\
and hollow Cauchy matrices}
\author{Principia Math in Progress\\MathDB problem 351647}
\date{29 September 2026}
\begin{document}
\maketitle
\begin{abstract}
We prove the two restricted reciprocal-permutation congruences for every
odd prime and every parameter prime to it, and the quadratic-residue
derangement congruence. For parameters divisible by the prime, the signed
formula holds for all primes at least five; at three its exact correction
is $-a$ modulo nine. Thus the literal conjecture has exactly the two
exceptional classes $p=3$, $a\equiv3,6\pmod9$. The proof uses a first-order
deformation of a hollow Cauchy matrix.
\end{abstract}

\paragraph{Source and scope.}
The target is Sun's Conjecture 4.4 in \cite{Sun}, page 26, reproduced in
\href{https://mathdb.com/p/351647}{MathDB 351647}. All congruences between
rationals are in $\Zp$, the rationals whose reduced denominators are
prime to $p$. The source allows every integer $a$ in the signed formula.
We retain that quantifier and identify all exceptions; the result below
also proves the remaining, unrestricted prime ranges.

\begin{theorem}\label{main}
Let $p$ be an odd prime, $n=p-1$, and $a\in\mathbb Z$. Define
\[
 M_p(a)_{ij}=
 \begin{cases}(a+ij)^{-1},&p\nmid a+ij,\\0,&p\mid a+ij,\end{cases}
 \qquad 1\leq i,j\leq n .
\]
If $p\nmid a$, then
\begin{align}
 \det M_p(a)&\equiv
 \left(\frac ap\right)\frac{3-a^{p-1}}2\pmod{p^2},\label{detmain}\\
 \per M_p(a)&\equiv
 (-1)^{(p+1)/2}\frac{3-a^{p-1}}2\pmod{p^2}.\label{permain}
\end{align}
If $p\mid a$, the signed congruence \eqref{detmain} holds when $p\geq5$.
For $p=3$ and $3\mid a$, instead
\begin{equation}
 \det M_3(a)\equiv-a\pmod9.\label{exception}
\end{equation}
Finally, if $p\equiv1\pmod4$ and $m=(p-1)/2$, then
\begin{equation}
 \sum_{\tau\in D(m)}
 \prod_{j=1}^m\frac1{j^2-\tau(j)^2}
 \equiv\big((m/2)!\big)^{-2}\pmod p.\label{quadratic}
\end{equation}
Here $D(m)$ is the set of permutations without fixed points.
\end{theorem}

The determinant and permanent of $M_p(a)$ are exactly the signed and
unsigned restricted sums: setting a forbidden entry to zero deletes
precisely the forbidden permutations.

\paragraph{Prior ingredients.}
Cycle cancellation for reciprocal differences already appears in
Guo--Li--Tao--Wei \cite[Lemma 3.1]{GLTW}. The resulting
permanent--determinant identity, and the canonical determinant modulo
$p^2$, appear in Yang--Zhang \cite[Proposition 12 and Lemma 15]{YZ}.
We give self-contained proofs and claim no separate problem credit for
these ingredients. The quadratic-residue calculation is also closely
related to the root-of-unity evaluations in \cite{GLTW}.

\begin{lemma}\label{cycles}
For an even integer $N$ and pairwise distinct nodes $x_1,\ldots,x_N$,
let $C(x)$ have off-diagonal entries $(x_i-x_j)^{-1}$ and diagonal zero.
As a rational identity,
\[
 \per C(x)=(-1)^{N/2}\det C(x).
\]
\end{lemma}
\begin{proof}
Fix a set of $\ell\geq3$ nodes and distinguish one of them, $z$.
Every oriented cycle on the set is obtained uniquely by inserting $z$
in one edge $(u,v)$ of a cycle on the other $\ell-1$ nodes. The factor
by which the cycle product changes is
\[
 \frac{u-v}{(u-z)(z-v)}
 =\frac1{u-z}-\frac1{v-z}.
\]
Summing this over the edges of any fixed contracted cycle gives zero.
Thus the sum of products over all cycles on a given set of at least
three nodes is zero. In either permutation expansion, group permutations
by the sets supporting their cycles. Fixed points contribute zero;
a block of size at least three cancels. The remaining permutations
are perfect matchings, all with sign $(-1)^{N/2}$. This proves the identity.
It remains valid over any field whenever the node differences are units.
\end{proof}

\begin{lemma}\label{canonical}
For $C_p=C(1,2,\ldots,p-1)$, one has
$\det C_p\equiv1\pmod{p^2}$.
\end{lemma}
\begin{proof}
Write $n=p-1$, $P(X)=\prod_{j=1}^n(X-j)$, $w_j=1/P'(j)$,
$W=\operatorname{diag}(w_j)$, and $X=\operatorname{diag}(1,\ldots,n)$.
The matrix $D$ that sends the values of a polynomial of degree below
$n$ to the values of its derivative has, by Lagrange interpolation,
\[
 D_{ij}=\frac{w_j}{w_i(i-j)}\quad(i\ne j),\qquad
 D_{ii}=\sum_{j\ne i}\frac1{i-j}=H_{i-1}-H_{n-i},
\]
where $H_0=0$ and $H_r=\sum_{j=1}^r1/j$.
Set $L=D+X^{-1}$. The matrix $XL=XD+I$ represents
$f\mapsto xf'+f$, whose eigenvalues on $1,x,\ldots,x^{n-1}$
are $1,\ldots,n$. Therefore
\[
 \det L=\frac{n!}{\det X}=1.
\]
Conjugation by $W$ shows that $L$ is similar to
$C_p+\operatorname{diag}(u_i)$, where
$u_i=H_i-H_{p-1-i}$. Since
$H_{p-1}\equiv0\pmod p$ and
$\sum_{j=p-i}^{p-1}1/j\equiv-H_i\pmod p$, every $u_i$ is in $p\Zp$.
In expanding $\det(C_p+\operatorname{diag}(u_i))$, each term linear
in a $u_i$ has an odd-order skew-symmetric principal minor as coefficient,
so that coefficient is zero. Terms containing at least two $u_i$ lie
in $p^2\Zp$. Consequently $\det C_p\equiv\det L=1\pmod{p^2}$.
\end{proof}

\begin{lemma}\label{deformation}
Let $x_i=i$ for $1\leq i\leq p-1$, and let $y_j=j+\delta_j$,
where $\delta_j\in p\Zp$. Define $H(x,y)$ to have zero diagonal and
off-diagonal entries $(x_i-y_j)^{-1}$. Then, with
$\chi=(-1)^{(p-1)/2}$,
\begin{align}
 \det H(x,y)&\equiv1-\frac12\sum_{j=1}^{p-1}\frac{\delta_j}{j}
                    \pmod{p^2},\label{detdeform}\\
 \per H(x,y)&\equiv\chi\det H(x,y)\pmod{p^2}.\label{perdeform}
\end{align}
\end{lemma}
\begin{proof}
Put $F(x,y)=\det H(x,y)$ and $G(x)=F(x,x)=\det C(x)$.
Since $H(y,x)=-H(x,y)^T$ and $n=p-1$ is even,
$F(x,y)=F(y,x)$. On the diagonal this gives
\[
 F_{y_j}(x,x)=F_{x_j}(x,x)=\tfrac12G_{x_j}(x).
\]
We compute the last derivative at $r=(1,\ldots,p-1)$ over $\Fp$.
The rational function $G$ is symmetric and homogeneous of degree $-n$.
For $g\in\Fp^\times$, multiplication by $g$ permutes the coordinates
of $r$. If $\pi_g(j)$ is the index whose residue is $gj$, symmetry and
homogeneity yield
\[
 G_{x_{\pi_g(j)}}(r)=g^{-n-1}G_{x_j}(r),\qquad
 r_{\pi_g(j)}G_{x_{\pi_g(j)}}(r)
     =g^{-n}r_jG_{x_j}(r)=r_jG_{x_j}(r).
\]
Thus all $r_jG_{x_j}(r)$ have the same value. Euler's identity says their
sum is $-nG(r)$. Since $n$ is nonzero in $\Fp$, their common value is
$-G(r)=-1$, by Lemma \ref{canonical}. It follows that
$F_{y_j}(r,r)=-1/(2j)$ in $\Fp$.
The constant term is $F(r,r)\equiv1\pmod{p^2}$. Taylor expansion in
the $\delta_j$ now proves \eqref{detdeform}.

For completeness set $R(x,y)=\per H(x,y)-\chi\det H(x,y)$.
It too is symmetric in $x,y$, and $R(x,x)=0$ identically by
Lemma \ref{cycles}. Differentiation gives
$R_{x_j}(x,x)+R_{y_j}(x,x)=0$, while symmetry makes the two terms equal.
They therefore both vanish. Its expansion at $(r,r)$ has neither
constant nor linear terms, proving \eqref{perdeform}.
All denominators in these expansions are products of off-diagonal
node differences, which are units modulo $p$; all terms of degree at
least two are consequently in $p^2\Zp$.
\end{proof}

\paragraph{Proof of the unit-parameter formulas.}
Assume $p\nmid a$. Let $\sigma$ be the permutation of $1,\ldots,n$
satisfying $\sigma(j)\equiv-a/j\pmod p$, and put
\[
 y_j=-\frac a{\sigma(j)},\qquad \delta_j=y_j-j,\qquad
 B=(p-1)!,\qquad A=a^{p-1}.
\]
Then $\delta_j\in p\Zp$. For the column-permuted matrix,
\[
 M_p(a)_{i,\sigma(j)}=\frac1{\sigma(j)}H(x,y)_{ij}.
\]
This equality also holds on the diagonal: both entries are forced to
zero there. It follows that
\begin{equation}
 \det M_p(a)=\frac{\sgn\sigma}{B}\det H(x,y),\qquad
 \per M_p(a)=\frac1B\per H(x,y).\label{column}
\end{equation}
The exact product identity
\[
 \prod_{j=1}^n\frac{y_j}{j}=\frac{A}{B^2}
\]
implies $\sum_j\delta_j/j\equiv A/B^2-1\pmod{p^2}$, so
\[
 \det H(x,y)\equiv\frac{3-A/B^2}{2}\pmod{p^2}.
\]
Write $B=-1+pw$ and $A=1+pq$, using Wilson's and Fermat's theorems.
Then
\[
 B^{-1}\frac{3-A/B^2}{2}
 \equiv(-1-pw)\left(1-pw-\frac{pq}{2}\right)
 \equiv-1+\frac{pq}{2}
 =-\frac{3-A}{2}\pmod{p^2}.\label{cancel}
\]
To evaluate $\sgn\sigma$, inversion on $\Fp^\times$ fixes $1,-1$
and has $(p-3)/2$ transpositions, hence sign $-\chi$.
Multiplication by $b\in\Fp^\times$ has sign $(b/p)$:
the ratio of Vandermonde products gives its sign modulo $p$ as
$b^{n(n-1)/2}=(b^{(p-1)/2})^{p-2}=(b/p)$, and the signs
$1,-1$ are distinct modulo an odd prime.
Taking $b=-a$ gives
$\sgn\sigma=(-\chi)(-a/p)=-(a/p)$.
Substitution in \eqref{column}, together with
Lemma \ref{deformation}, proves \eqref{detmain} and \eqref{permain}.

\paragraph{Parameters divisible by the prime.}
If $p\mid a$, every entry is allowed and
$M_p(a)_{ij}\equiv1/(ij)\pmod p$.
For each $i>1$, subtract $i^{-1}$ times the first row from row $i$.
All entries in those $n-1$ rows now lie in $p\Zp$, so
$\det M_p(a)\in p^{n-1}\Zp$. For $p\geq5$, this is zero modulo $p^2$,
as is the right side of \eqref{detmain}.
For $p=3$ and $3\mid a$, direct calculation gives
\[
 \det M_3(a)=
 \frac1{(a+1)(a+4)}-\frac1{(a+2)^2}
 =\frac{-a}{(a+1)(a+4)(a+2)^2}
 \equiv-a\pmod9.
\]
The denominator is $1$ modulo $3$, and the numerator is divisible by
$3$, which justifies the last reduction modulo $9$.
Thus the original zero prediction fails exactly for
$a\equiv3,6\pmod9$, and succeeds for $a\equiv0\pmod9$.

\paragraph{Proof of the quadratic-residue formula.}
Work over $\Fp$. Since $p\equiv1\pmod4$, the integer $m=(p-1)/2$
is even. The nodes $x_j=j^2$, $1\leq j\leq m$, are distinct and are
precisely the roots of $P(T)=T^m-1$. Let $Q=C(x)$ and
$X=\operatorname{diag}(x_j)$.
The same Lagrange differentiation matrix as in Lemma \ref{canonical}
has diagonal
\[
 D_{jj}=\frac{P''(x_j)}{2P'(x_j)}
       =\frac{m-1}{2x_j}.
\]
Writing $c=(m-1)/2$ in $\Fp$, diagonal conjugation therefore makes
$Q$ similar to $D-cX^{-1}$. On polynomials of degree below $m$,
$XD$ represents $f\mapsto xf'$, with eigenvalues $0,1,\ldots,m-1$.
Consequently
\[
 \det Q=\frac{\prod_{d=0}^{m-1}(d-c)}{\prod_{j=1}^m x_j}
 =\frac{(-1)^{m/2}((m-1)!!)^2}{2^m(m!)^2}.
\]
By Lemma \ref{cycles} and $m!=2^{m/2}(m/2)!(m-1)!!$, this gives
\[
 \per Q=\frac{1}{2^{2m}((m/2)!)^2}
       =\frac1{((m/2)!)^2}\quad\hbox{in }\Fp,
\]
where $2^{2m}=2^{p-1}=1$. The zero diagonal makes $\per Q$ the sum
in \eqref{quadratic}, completing Theorem \ref{main}.
\qed

\paragraph{Verification and limits.}
The accompanying standard-library Python verifier evaluates the original
matrices using fraction-free integer determinants and subset dynamic
programming for permanents. It checks every parameter modulo $p^2$
for $p=3,5,7,11,13$, additional larger-prime determinants, the
quadratic-residue formula, and arbitrary first-order column deformations.
Exact rational permutation enumeration supplies a separate small-matrix
cross-check. These are corroborating computations; the preceding proof
establishes the infinite ranges. The package has author-agent review,
not external peer review.

\paragraph{Task 2. A lift modulo $p^3$.}
Determine the second-order deformation of $H(x,y)$ and the first
possible correction to \eqref{detmain}--\eqref{permain} modulo $p^3$.
Express the answer using node displacements and harmonic sums.

\paragraph{Task 3. Changing integer representatives.}
Replace $1,\ldots,p-1$ by arbitrary integer representatives of
$\Fp^\times$ in the row and column sets. Determine exactly which
linear displacement statistics govern the resulting sums modulo $p^2$.

\paragraph{Task 4. Other multiplicative subgroups.}
For each even divisor $m$ of $p-1$, and a coset of the subgroup of
order $m$, evaluate the hollow reciprocal-difference permanent.
Determine which such evaluations lift from $\Fp$ to $\mathbb Z/p^2\mathbb Z$.

\paragraph{Task 5. More than one forbidden residue.}
Replace the denominator $a+ij$ by a polynomial with two or more
forbidden residue classes. Identify conditions under which a
permanent--determinant relation survives to first order.

\begin{thebibliography}{9}
\bibitem{Sun} Z.-W. Sun, \emph{Arithmetic properties of some permanents},
\href{https://arxiv.org/abs/2108.07723v7}{arXiv:2108.07723v7} (2022),
Conjecture 4.4, p.~26.
\bibitem{GLTW} X. Guo, X. Li, Z. Tao and T. Wei,
\emph{The eigenvectors-eigenvalues identity and Sun's conjectures on
determinants and permanents},
\href{https://arxiv.org/abs/2206.02592}{arXiv:2206.02592} (2022);
\emph{Linear Multilinear Algebra} \textbf{72} (2024), 1071--1077,
\href{https://doi.org/10.1080/03081087.2023.2172380}{doi:10.1080/03081087.2023.2172380}.
The lemma locator above refers to the arXiv version.
\bibitem{YZ} Y. Yang and Y. Zhang,
\emph{Sun-type determinant and permanent congruences},
\href{https://arxiv.org/abs/2605.19502v3}{arXiv:2605.19502v3} (2026).
\end{thebibliography}
\end{document}
